\documentclass[11pt,a4paper]{article}
\usepackage[T1]{fontenc}
\usepackage[utf8]{inputenc}
\usepackage{lmodern,amsmath,amssymb}
\usepackage[margin=25mm]{geometry}
\usepackage{longtable,booktabs,array,calc}
\usepackage[hidelinks]{hyperref}
\hypersetup{pdftitle={The record's distance is a path length},pdfauthor={navstokgap research working notes}}
\newcommand{\tightlist}{\setlength{\itemsep}{0pt}\setlength{\parskip}{0pt}}
\setlength{\emergencystretch}{3em}
\setcounter{secnumdepth}{0}
\begin{document}
\hypertarget{the-records-distance-is-a-path-length}{%
\section{The record's distance is a path
length}\label{the-records-distance-is-a-path-length}}

Three results that close STATE's next items 1 and 2 of 2026-09-23.

\textbf{The single-shot constant for every protocol.} Every protocol
theorem of the programme --- the aperture bound of the
\href{planck-gap-probabilistic.md}{probabilistic note}, the recoil bound
of the \href{record-costs-recoil.md}{recoil note}, the disturbance bound
of the \href{record-costs-disturbance.md}{disturbance note} --- holds
with \(\arcsin(1-2\epsilon)\) in place of \(1-2\epsilon\) (Theorem P).
Since \(\arcsin(1-2\epsilon)=\arccos(2\sqrt{\epsilon(1-\epsilon)})\),
this is the sharp single-shot constant of Theorem 1 of the probabilistic
note, so the gap between the single-shot and protocol forms is closed,
and at certain decision the recoil bound reads

\[s\sum_j\Delta_j\ \ge\ 4\pi\hbar=2h .\]

\textbf{The aperture bound in its kick form.} For every protocol of any
instruments, with the same preparation under both hypotheses,

\[\hbar\arcsin(1-2\epsilon)\ \le\ \int_0^\tau|f(t)|\,L(t)\,dt ,\]

with \(L(t)\) the body's position spread at time \(t\) (Theorem K). A
force is a succession of impulses, and an impulse is registered only
through the position spread it acts on. The \((L,P)\) form of the
probabilistic note follows because free motion spreads position at the
rate \(P/m\), and the two-packet preparation saturates the kick form
exactly as \(\epsilon\to0\).

\textbf{One accounting.} The aperture bound pays for the force at the
body, the disturbance bound pays for it at the marks. Any split of the
force between the two gives a valid bound (Theorem M), and for von
Neumann marks and a force of one sign the best split is pointwise:

\[\hbar\arcsin(1-2\epsilon)\ \le\ \int_0^\tau|f(u)|\,\min\bigl(L(u),R(u)\bigr)\,du,\qquad
R(u)=\frac1m\sum_j\Delta_j\,G_\tau(t_j,u),\]

with \(G_\tau\) the Dirichlet Green's function of the interval. \(R(u)\)
is the \textbf{recoil length}: the marks' recoil spreads weighted by the
displacement each one would produce at time \(u\) relative to the chord
pinned at the ends, summed, which is a length (an upper bound on the
spread those recoils produce, since independent recoils add in
quadrature). So the force is paid in length, at every moment the smaller
of the body's spread and the marks' recoil length (Corollary M1).

\emph{Prior art.} Lemma 1 is the mixed-state Mandelstam--Tamm bound in
the Bures angle (Uhlmann,
\href{https://doi.org/10.1016/0375-9601(92)90555-Z}{Phys. Lett. A
\textbf{161}, 329, 1992}; Anandan and Aharonov,
\href{https://doi.org/10.1103/PhysRevLett.65.1697}{PRL \textbf{65},
1697, 1990}); statistical distance as a path length is Wootters's
(\href{https://doi.org/10.1103/PhysRevD.23.357}{PRD \textbf{23}, 357,
1981}) and Braunstein and Caves's
(\href{https://doi.org/10.1103/PhysRevLett.72.3439}{PRL \textbf{72},
3439, 1994}); the hybrid chain is the hybrid argument of Bennett,
Bernstein, Brassard and Vazirani
(\href{https://doi.org/10.1137/S0097539796300933}{SIAM J. Comput.
\textbf{26}, 1510, 1997}); Bures-angle speed limits for general
processes are those of Taddei and coauthors
(\href{https://doi.org/10.1103/PhysRevLett.110.050402}{PRL \textbf{110},
050402, 2013}); and the position-spread form of force sensitivity is the
waveform estimation limit of Tsang, Wiseman and Caves
(\href{https://doi.org/10.1103/PhysRevLett.106.090401}{PRL \textbf{106},
090401, 2011}); all metadata. New here are the split accounting of
Theorem M and the recoil length of Corollary M1. Exploratory; no ledger
promotion.

\hypertarget{the-bures-angle}{%
\subsection{1. The Bures angle}\label{the-bures-angle}}

For density operators \(\rho,\sigma\) let
\(F(\rho,\sigma)=\operatorname{Tr}|\sqrt\rho\sqrt\sigma|\) be the root
fidelity and \(A(\rho,\sigma)=\arccos F(\rho,\sigma)\in[0,\pi/2]\) the
Bures angle. Four standard facts are used.

\begin{enumerate}
\def\labelenumi{\arabic{enumi}.}
\tightlist
\item
  \(A\) is a metric on states: it obeys the triangle inequality.
\item
  \(A\) does not increase under any channel, measurements and classical
  post-processing included (monotonicity of the fidelity).
\item
  \(F(\rho,\sigma)=\max|\langle\psi|\phi\rangle|\) over purifications
  (Uhlmann,
  \href{https://doi.org/10.1016/0034-4877(76)90060-4}{Rep.~Math. Phys.
  \textbf{9}, 273, 1976}, metadata), so \(A(\rho,\sigma)\) is at most
  the Fubini--Study angle between any two purifications.
\item
  \(\tfrac12\|\rho-\sigma\|_1\le\sin A(\rho,\sigma)\) (Fuchs and van de
  Graaf, \href{https://doi.org/10.1109/18.761271}{IEEE Trans. Inf.
  Theory \textbf{45}, 1216, 1999}, metadata); for probability
  distributions this is the classical statement.
\end{enumerate}

\textbf{Lemma 1 (one unitary step).} Let \(\rho\) be a state of body,
apparatus and a classical record register \(x\),
\(\rho=\sum_xp_x|x\rangle\langle x|\otimes\rho_x\), and let
\(V=\sum_x|x\rangle\langle x|\otimes V_x\) with
\(V_x=\mathcal T\exp(i\int_0^1K_x(u)\,du)\). Then
\(A(\rho,V\rho V^\dagger)\le\sup_x\int_0^1\Delta K_x(u)\,du\), each
spread taken in the state \(V_x(u)\rho_xV_x(u)^\dagger\) along the path.

\emph{Proof.} Square roots of block-diagonal operators are block
diagonal, so
\(F(\rho,V\rho V^\dagger)=\sum_xp_xF(\rho_x,V_x\rho_xV_x^\dagger)\), and
\(\sum_xp_x\cos A_x\ge\cos\sup_xA_x\). For each block, purify \(\rho_x\)
to \(\psi_x\); by fact 3 the angle is at most the Fubini--Study length
of the curve \(V_x(u)\psi_x\), whose speed is the spread of \(K_x(u)\)
in the current state, and that spread equals the spread in the mixed
state it purifies. \(\square\)

For a time-independent generator \(G\) acting on the body, the bound is
\(\sup_x\Delta_{\rho_x}G\), the spread in the body's marginal.

\hypertarget{theorem-p-the-single-shot-constant-for-every-protocol}{%
\subsection{2. Theorem P: the single-shot constant for every
protocol}\label{theorem-p-the-single-shot-constant-for-every-protocol}}

\textbf{Theorem P.} In Theorem 2 of the probabilistic note, Theorem R
and Corollary R1 of the recoil note, Theorem U and Corollary U1 of the
disturbance note, and Theorems K and M below, the conclusion holds with
\(\arcsin(1-2\epsilon)\) in place of \(1-2\epsilon\).

\emph{Proof.} Each proof builds a chain of hybrid processes
\(H_0,\dots,H_n\) joining the two hypotheses, in which consecutive
hybrids apply the same channel before and after one step where they
differ by a unitary. Replace the triangle inequality for total variation
by fact 1, data processing by fact 2, and the per-step trace-distance
bound by Lemma 1. This gives
\(A(P_{\mathrm I},P_{\mathrm F})\le\Sigma\), the same sum \(\Sigma\) of
spreads over \(\hbar\) that each proof obtains. A test with equal-prior
error at most \(\epsilon\) has
\(\tfrac12\|P_{\mathrm I}-P_{\mathrm F}\|_1\ge1-2\epsilon\), so by fact
4 \(\sin A\ge1-2\epsilon\), hence
\(\Sigma\ge A\ge\arcsin(1-2\epsilon)\). In the recoil theorem the two
probe families are products, and the chain replaces one factor at a
time. \(\square\)

The constant is sharp in the single-shot and aperture forms. For a
single shot it is Theorem 1 of the probabilistic note, the aperture
bound's protocol form is now the same inequality with the same constant,
and the two-packet preparation attains the kick form as \(\epsilon\to0\)
(§3). The recoil and disturbance bounds are not shown attained: the best
known recoil protocol sits a factor of about \(1.47\) above the bound at
five per cent error. At \(\epsilon\to0\) the constant is \(\pi/2\) where
it was \(1\).

\hypertarget{theorem-k-the-aperture-bound-in-its-kick-form}{%
\subsection{3. Theorem K: the aperture bound in its kick
form}\label{theorem-k-the-aperture-bound-in-its-kick-form}}

Write \(f(t)\) for the force history, \(F\)-hypothesis against free
motion, with the same preparation \(\sigma\) under both. Let \(L(t)\) be
the supremum of the body's position spread \(\Delta\hat y\) at time
\(t\), in the laboratory frame, over the conditional states that can
occur at time \(t\) when the force acts on any initial part of the
interval.

\textbf{Theorem K.} Every protocol of instruments, of any kind, adaptive
or not, deciding at error \(\epsilon\) satisfies
\(\hbar\arcsin(1-2\epsilon)\le\int_0^\tau|f(t)|L(t)\,dt\).

\emph{Proof.} Let hybrid \(H_{t^*}\) apply the force on \([0,t^*]\) and
none afterwards, the marks unchanged; \(H_\tau\) is \(\mathrm F\) and
\(H_0\) is \(\mathrm I\). Take a partition containing the mark times. On
a cell \([t,t+\delta]\) free of marks, \(H_{t+\delta}\) and \(H_t\)
agree up to \(t\) and after \(t+\delta\), and differ by the forced
evolution over the cell, which for linear dynamics is the free evolution
followed by the translation by the displacement the force produces from
rest,
\((\int_t^{t+\delta}(t+\delta-s)f\,ds/m,\ \int_t^{t+\delta}f\,ds)\). Its
generator has spread at most
\(|\int_t^{t+\delta}f|\,\Delta\hat y+O(\delta^2)\Delta\hat p\) in the
state it acts on. Lemma 1 and Theorem P's chain give
\(\hbar A\le\sum_{\rm cells}|\int f|\,L+O(\delta)\sup\Delta\hat p\), and
refining the partition gives the integral when the momentum spread of
the chain's states is bounded. If it is infinite for some conditional
state on some cell, free evolution makes that state's position spread
infinite at every later time of the cell, so \(L=\infty\) on a set of
positive measure and the theorem holds vacuously. \(\square\)

\textbf{Corollary K1 (the aperture form).} If the conditional states
after each mark, and at \(t=0\), have position spread at most \(L\) and
momentum spread at most \(P\), then between marks the spread grows at
most linearly, \(L(t)\le L+P(t-t_j)/m\), because a c-number force
translates and free motion shears. For the constant force,

\[F\tau L+\frac{F\tau^2}{2m}P\ \ge\ \hbar\arcsin(1-2\epsilon),\]

and with marks every \(\delta t\) the second term shrinks to at most
\(F\tau\,\delta t\,P/2m\). \(\square\)

This is Theorem 2 of the probabilistic note with the sharp constant, and
it holds under either reading of that note's aperture. With spreads
bounded in the interaction picture, \(\Delta(\hat y-\hat pt/m)\le L\),
the laboratory spread is at most \(L+Pt/m\), and the same integral
results.

\textbf{Saturation.} The two-packet preparation of the probabilistic
note, separation \(\pi\hbar/(F\tau)\), has position spread
\(\pi\hbar/(2F\tau)\) up to a relative correction of order
\(\sqrt{\tau\Delta E/\hbar}\) from spreading. So
\(\int_0^\tau F\,L\,dt\to\pi\hbar/2\) as \(\tau\Delta E\to0\), while its
error probability tends to zero and \(\arcsin(1-2\epsilon)\to\pi/2\).
The kick form is attained exactly in that limit.

\hypertarget{theorem-m-one-accounting}{%
\subsection{4. Theorem M: one
accounting}\label{theorem-m-one-accounting}}

\textbf{Theorem M.} Split the force as \(f=g+h\), and let \(c\) be any
function on \([0,\tau]\) with \(c''=g/m\). For every protocol of
instruments deciding at error \(\epsilon\) whatever the body's initial
state under each hypothesis (the composite requirement of Theorem U;
with the same initial state assumed under both, the chain's mixed pair
would only give \(\arcsin(1-4\epsilon)\)),

\[\hbar\arcsin(1-2\epsilon)\ \le\ \int_0^\tau|h(t)|\,L(t)\,dt
+\sum_j\sup\Delta\bigl(c(t_j)\hat D_j-mc'(t_j)\hat X_j\bigr),\]

with \(\hat D_j,\hat X_j\) the disturbances of mark \(j\) as in the
disturbance note, and \(L(t)\) taken over the states of the chain below.

\emph{Proof.} Take the pair of initial states \(\sigma\) under
\(\mathrm I\) and \(T_c(0)\sigma\) under \(\mathrm F\), \(T_c(t)\) the
translation by \((c(t),mc'(t))\). The chain has two legs. First,
starting from the process with force \(f\) from \(T_c(0)\sigma\), switch
off \(h\) cell by cell as in Theorem K, paying \(\int|h|L/\hbar\); this
ends at the process with force \(g\) from \(T_c(0)\sigma\). Since
\(c''=g/m\), the process with force \(g\) from \(T_c(0)\sigma\) is the
free process from \(\sigma\) translated by \(T_c(t)\) between marks,
which is the setting of Theorem U; its chain reaches the free process
from \(\sigma\) and pays the second sum. Lemma 1 bounds every step and
fact 1 adds the legs. \(\square\)

With \(h=f\) and \(c=0\) it is Theorem K; with \(h=0\) and \(c=P-a-bt\)
it is Theorem U. The theorem reads as a transport of the offset between
the hypotheses. A free line of offset is absorbed in the initial state
at no cost, curvature carried through a mark costs that mark's
disturbance weighted by the offset it meets, and force removed at the
body costs impulse times position spread. The record's Bures distance is
at most the length of any such transport, and a protocol escapes the
floor only if at every moment the body is spread, or the marks disturb,
enough to cover the force acting then.

\textbf{Corollary M1 (the force is paid in length).} For von Neumann
marks with recoil spreads \(\Delta_j\) and a force of one sign, take
\(g=\chi f\) with \(0\le\chi\le1\) and \(c\) the solution vanishing at
\(0\) and \(\tau\), \(c(t)=-\frac1m\int_0^\tau G_\tau(t,u)g(u)\,du\),
\(G_\tau(t,u)=\min(t,u)(\tau-\max(t,u))/\tau\). Then
\(\sum_j\Delta_j|c(t_j)|=\int\chi|f|R\), and minimizing over \(\chi\)
pointwise, which is legitimate because for von Neumann marks the chain's
conditional spreads do not depend on \(\chi\),

\[\hbar\arcsin(1-2\epsilon)\ \le\ \int_0^\tau|f(u)|\min\bigl(L(u),R(u)\bigr)\,du,
\qquad R(u)=\frac1m\sum_j\Delta_jG_\tau(t_j,u). \qquad\square\]

A mark at \(0\) or at \(\tau\) contributes nothing to \(R\), whatever
its recoil, since \(G_\tau\) vanishes there; its reading still fixes the
chord. For the three marks of Proposition D of the
\href{mark-cost-and-statistical-floor.md}{mark-cost note}, outer marks
sharp, \(R(u)=\Delta_2\min(u,\tau-u)/2m\), a tent of height
\(\Delta_2\tau/4m\). The pure recoil bound is \(F\int R=\Delta_2s/4\);
the pure aperture bound is \(F\int L\). The recoil length vanishes at
the ends and peaks in the middle, so the mixed bound pays at the marks
near the ends and at the body wherever the body is more compact than
\(R\). Whenever \(L\) and \(R\) cross, \(\int F\min(L,R)\) is strictly
below both pure costs, and the necessary condition it states is strictly
stronger than either.

\hypertarget{the-numbers}{%
\subsection{5. The numbers}\label{the-numbers}}

\begin{longtable}[]{@{}
  >{\raggedright\arraybackslash}p{(\columnwidth - 8\tabcolsep) * \real{0.2000}}
  >{\raggedright\arraybackslash}p{(\columnwidth - 8\tabcolsep) * \real{0.2000}}
  >{\raggedright\arraybackslash}p{(\columnwidth - 8\tabcolsep) * \real{0.2000}}
  >{\raggedright\arraybackslash}p{(\columnwidth - 8\tabcolsep) * \real{0.2000}}
  >{\raggedright\arraybackslash}p{(\columnwidth - 8\tabcolsep) * \real{0.2000}}@{}}
\toprule\noalign{}
\begin{minipage}[b]{\linewidth}\raggedright
Bound
\end{minipage} & \begin{minipage}[b]{\linewidth}\raggedright
Old constant
\end{minipage} & \begin{minipage}[b]{\linewidth}\raggedright
Sharp constant
\end{minipage} & \begin{minipage}[b]{\linewidth}\raggedright
At \(\epsilon=0.05\)
\end{minipage} & \begin{minipage}[b]{\linewidth}\raggedright
Certain decision
\end{minipage} \\
\midrule\noalign{}
\endhead
\bottomrule\noalign{}
\endlastfoot
Recoil, \(s\sum_j\Delta_j/\hbar\ge\) & \(8(1-2\epsilon)\) &
\(8\arcsin(1-2\epsilon)\) & \(8.96\) & \(4\pi\),
i.e.~\(s\sum\Delta_j\ge2h\) \\
Disturbance,
\(\bigl(\frac s8\sum\Delta(\hat D_j)+\frac J2\sum\Delta(\hat X_j)\bigr)/\hbar\ge\)
& \(1-2\epsilon\) & \(\arcsin(1-2\epsilon)\) & \(1.12\) & \(\pi/2\) \\
Aperture, \(\bigl(F\tau L+\frac{F\tau^2}{2m}P\bigr)/\hbar\ge\) &
\(1-2\epsilon\) & \(\arcsin(1-2\epsilon)\) & \(1.12\) & \(\pi/2\) \\
\end{longtable}

The balanced aperture of minimal area now gives
\(\tau\Delta E\ge\frac12\hbar\arcsin^2(1-2\epsilon)\), which is
\(\pi^2\hbar/8\) at certain decision. The three-mark protocol with the
best Gaussian probe sits within \(z_{1-\epsilon}/\arcsin(1-2\epsilon)\)
of the recoil bound, about \(1.47\) at five per cent error, at every
recoil (recoil note, §4).

\hypertarget{consequence-for-state}{%
\subsection{6. Consequence for STATE}\label{consequence-for-state}}

Items 1 and 2 of STATE's queue are settled, apart from the worst-case
interval \([9,36]\) of the derivation note, which concerns a framework
without quantum instances. Every protocol theorem carries the sharp
single-shot constant; the aperture bound is the kick form
\(\int|f|L\ge\hbar\arcsin(1-2\epsilon)\); and one accounting, Theorem M,
contains the aperture and disturbance bounds and interpolates between
them, with the pointwise form \(\int|f|\min(L,R)\) for von Neumann
marks. The remaining modern item is whether the minimum over splits and
lines in Theorem M is attained by some protocol, which would make it the
exact floor.
\end{document}
